\section{Two further coefficient-three labels from a certified real class group}
\label{sec:coefficient-three-conductor32-half}

\paragraph{The extra class input is explicit.}
The even primitive conductor-32 labels were excluded from G178 because
the two odd special values were not simultaneously units.
A certified real class group supplies the even component directly.
The finite computation below certifies a class-group quotient;
it is not a claim to have certified a complete unit basis.

\begin{theorem}[AC-G180: two conductor-32 rows on the original coefficient ring]
\label{AC-G180}
The real cyclotomic field $F=\mathbb Q(\zeta_{96}+\zeta_{96}^{-1})$
has class number one, as certified below.
Let $\delta$ be an even primitive Dirichlet character of conductor $32$.
Let $O$ be the original ring of integers of a finite extension
$E/\mathbb Q_3$ containing its values.  Write
$\beta=\overline{\delta(5)}$ in its residue field, and suppose
\[
                           \beta^2+\beta-1=0.
 \tag{G180.1}
\]
For every $e\geq1$, $s\in\mathbb Z$ and continuous
$\kappa:G_K\to1+\mathfrak m_O$ unramified away three,
\[
 \mathcal S_3\bigl(K,(E/O)
              (\delta\,\omega_3^{-s}\kappa^{-1}c_3^{-e})\bigr)=0.
 \tag{G180.2}
\]
The residual strict group is zero too.
For every actual fixed coefficient completion containing the required
eighth roots, exactly two of the four even primitive conductor-32
characters satisfy (G180.1).  They are distinct from every label of
G178.  No assertion is made here for the other two.
\end{theorem}

\begin{proof}
\proofstep{1. The field and exact class-group certificate.}
Define $S_0(X)=2$, $S_1(X)=X$, and
$S_{n+1}(X)=XS_n(X)-S_{n-1}(X)$.
Then $S_n(z+z^{-1})=z^n+z^{-n}$ by induction.
For a primitive $96$th root $z$ one has
$z^{16}+z^{-16}=1$.
Consequently $\alpha=z+z^{-1}$ is a root of
\[
 \begin{split}
 P(X)={}&S_{16}(X)-1\\
 ={}&X^{16}-16X^{14}+104X^{12}-352X^{10}
       +660X^8-672X^6+336X^4-64X^2+1.
 \end{split}\tag{G180.3}
\]
Conversely $\mathbb Q(z)$ is quadratic over its real subfield and
has degree $\varphi(96)=32$, so $\alpha$ has degree $16$.
This proves that $P$ is its minimal polynomial and identifies the
field used by the computation without an unchecked polynomial choice.

The retained exact PARI/GP script uses version 2.19.1,
$128$ decimal digits, and the explicit integer polynomial (G180.3).
It constructs $B=\operatorname{bnfinit}(P,1)$ and reports
\[
 \begin{array}{c|c}
 \operatorname{polisirreducible}(P)&1\\
 B.\mathrm{no}&1\\
 B.\mathrm{cyc}&[\,]\\
 B.\mathrm{disc}&121029087867608368152576\\
 \operatorname{bnfcertify}(B,1)&1.
 \end{array}\tag{G180.4}
\]
The preliminary \texttt{bnfinit} class output alone is conditional;
it is not the certificate used in this proof.
The final flag-one certification has the following precise documented
meaning: the actual class group is a quotient of the computed group.
Since the computed group is trivial, the actual class group is trivial.
The primary interface is the PARI/GP manual, section
\texttt{bnfcertify}, especially its flag-one clause, and section
\texttt{bnfinit} for the distinction from the preliminary computation.
The exact input, output and local primary-manual excerpt are retained.
No GRH assumption or unproved extension of theorem A is used to
replace this completed certificate.
The separate full-unit certification attempt is not needed and
produced no full certificate; its scope is recorded separately.

\proofstep{2. The even class sector of the full cyclotomic field.}
Put $H=\mathbb Q(\zeta_{96})$ and
$A=\operatorname{Cl}(H)\{3\}$.
Let $\iota$ be complex conjugation and let $N$ and $j$ denote
the class-group norm and ideal extension for $H/F$.
The exact identity is
\[
                         jN=1+\iota.
 \tag{G180.5}
\]
Step 1 makes the image of $N$ trivial.  On the even part $A^+$,
the right side of (G180.5) is multiplication by two.
That is an automorphism on a three-primary finite group.
Thus $A^+=0$, including the rational character component $\delta$.
Only class-number one, not a regulator or unit-basis claim, is used.

\proofstep{3. The complementary odd sector and the embedding criterion.}
The other rational extension of $\delta|_{G_K}$ is the odd
character $\nu=\delta\chi_{-3}$, of primitive conductor $96$.
Equation (G178.12) gives, for a primitive eighth root $\gamma$,
\[
 \overline{L(0,\delta\chi_{-3})}
       =(\gamma+1)(\gamma^2+\gamma-1),
                    \qquad\gamma=\overline{\delta(5)}.
 \tag{G180.6}
\]
To annihilate the $\nu$ class component, the required value is
$L(0,\nu^{-1})=L(0,\delta^{-1}\chi_{-3})$.
Thus the relevant argument of the right side is $\gamma=\beta^{-1}$.
Since $\beta^2+\beta-1=0$,
\[
   \beta^2-\beta-1=-2\beta=\beta\ne0
                                  \quad\text{in characteristic three}.
 \tag{G180.7}
\]
The two quadratics in
$X^4+1=(X^2+X-1)(X^2-X-1)$ are relatively prime,
and inversion exchanges their roots.
It follows that $\beta^{-1}+1$ and
$\beta^{-2}+\beta^{-1}-1$ are both nonzero.
The needed inverse special value is therefore a three-adic unit.
Its integrality follows either from the exact polynomial value
in (G178.12) or from the finite weighted sum before reduction.

Take $b_0\equiv5\pmod{32}$, $b_0\equiv1\pmod3$.
Then $b_0-\nu(b_0)$ reduces to $1-\beta$, which is nonzero.
The integral primitive Stickelberger annihilator of (G177.9)--(G177.10)
acts on $A(\nu)$ by the unit
$-(b_0-\nu(b_0))L(0,\nu^{-1})$.
Hence $A(\nu)=0$.
Together with Step 2 this kills both rational extensions of the
original residual character.  It does not assume that both
inverse special values are units.

\proofstep{4. Residual descent, full local condition and original lattice.}
The field $H$ contains $K$ with $[H:K]=16$, prime to three.
The restriction and class-field argument (G177.12)--(G177.13)
injects the strict residual group into
$\operatorname{Hom}_{\mathbb F_3[\operatorname{Gal}(H/K)]}
(A/3A,k'(\overline\delta))$.
Both rational components just proved zero imply this Hom is zero.
Faithful finite residue-field extension preserves the conclusion
over the original residue field.

At $K_3$, the finite character is unramified and a Frobenius lift
has value $\delta(3)$ of exact order eight.
The residue degree of $K_3/\mathbb Q_3$ is one.
The cyclotomic and $\kappa$ reductions on $G_K$ are one,
so the original scalar minus one is an $O$-unit at that lift.
The full local and global invariants of the discrete module vanish.
Away three, finite inertia has two-power order; its invariants
are either the full divisible module or zero.
The original-uniformizer exact sequence used in G178 Step 4 therefore
lifts every strict class of order $\varpi$ to the zero strict
residual group.  Every continuous cocycle is killed by a power of
$\varpi$, proving (G180.2) for the original, arbitrarily ramified $O$.

\proofstep{5. Exactly two actual characters.}
An even primitive character modulo $32$ is determined by its
value at $5$, which can be any primitive eighth root.
Reduction at three preserves this order.
The quadratic $X^2+X-1$ has two distinct roots in $\mathbb F_9$,
both primitive eighth roots: it is one of the two relatively prime
factors of $X^4+1$, and its discriminant is $2\ne0$ in $\mathbb F_3$.
Thus exactly two of the four actual characters satisfy (G180.1).
The condition is preserved by residue-field Frobenius and remains
well-defined after further original coefficient ramification.
For $h\geq3$ these characters occur in the actual host and their
values already belong to its original $O_3$.
The canonical identification (OPW.7) and the same inflation argument
therefore give the actual strict zeros at all $1\leq e<T$.
They have even parity at conductor $32$, whereas G178 contains
only the odd primitive characters at that conductor, so there
is no overlap.
\end{proof}

\paragraph{Boundary and geometric dictionary.}
The result adds two original labels at coefficient three.
It supplies no numerical all-prime bound or new mixed-cubic label.
The original finite-cover, Kummer and Tate-cohomology dictionaries
remain those of G177--G178, with their exact SGA 1 and Stacks scopes.
The class computation is an exact algorithmic certificate at the
documented PARI interface, not a Lean proof or full unit certificate.
The numerical remainder, $\mathcal R_{\rm cub}$ and H04.4 remain open.
